\section{Conclusiones}

\subsection{Síntesis por serie}

El objetivo planteado por la cátedra --determinar si los modelos contemporáneos de Machine Learning, Deep
Learning, modelos híbridos, AutoML y modelos de fundación mejoran el pronóstico obtenido con el enfoque
clásico del Trabajo Práctico N.\textsuperscript{o} 1-- admite una respuesta distinta según la serie y según
el criterio utilizado. Para la serie ALB, el modelo seleccionado mediante la regla honesta de validación
(AutoARIMA) mejora al SARIMA refit en un 48,5\,\% durante la validación, pero resulta un 34,9\,\% peor que
ese mismo SARIMA en el conjunto de prueba; de las 32 configuraciones evaluadas, el propio SARIMA refit del
trabajo anterior es, en definitiva, la de mejor desempeño de prueba. En esta serie, entonces, el enfoque
clásico del trabajo anterior no fue superado por ningún modelo más reciente cuando se lo evalúa sobre datos
no vistos durante el ajuste. Para la serie store-service, el modelo seleccionado (la referencia estacional
ingenua) tampoco es superado en validación por ningún modelo no contaminado, pero en el conjunto de prueba
sí lo supera un modelo de la familia Prophet (NeuralProphet, un 25,8\,\% mejor MASE que el modelo
seleccionado), aunque ese modelo no habría sido elegido por la regla honesta de selección, dado su
desempeño mediocre en validación (decimocuarto puesto de 32). La respuesta general, entonces, es matizada:
los modelos más recientes pueden superar al enfoque clásico, pero no necesariamente lo hacen los mismos
modelos que una regla de selección basada exclusivamente en validación elegiría.

\subsection{Consistencia entre validación y prueba}

La correlación de Spearman entre el orden de validación y el de prueba resultó baja y no significativa
para ALB ($\rho=0,237$, $p=0,192$ sobre las 32 configuraciones) y moderada y significativa para
store-service ($\rho=0,654$, $p<0,001$). Esta diferencia entre series sugiere que la propia confiabilidad
del protocolo de validación --y, por lo tanto, la confianza que puede depositarse en el modelo
seleccionado-- no es uniforme: para ALB, una ventana de validación de 48 horas ordena a los modelos de un
modo que guarda poca relación con su desempeño posterior, mientras que para store-service ese orden es más
informativo, aunque tampoco infalible, como lo demuestra el caso de NeuralProphet.

\subsection{Serie pública de contraste}

Sobre la serie de visitas horarias a Wikipedia en español, con seis modelos de hiperparámetros por defecto,
la regla de selección por validación eligió la referencia estacional ingenua, mientras que en prueba el mejor
modelo fue Ridge y el SARIMA del trabajo anterior el peor (2,513). El resultado es
coherente con la conclusión general del trabajo: el desempeño relativo del enfoque clásico depende de la
serie (fue el mejor en ALB y el peor en Wikipedia), y el orden de validación no anticipa el de prueba (Spearman
de 0,09 sobre seis modelos). Por su alcance reducido, este resultado se interpreta como un contraste y no
como una evaluación completa.

\subsection{Límites del trabajo}

\begin{enumerate}
  \item \textbf{Ventana única de validación.} El protocolo utiliza una sola ventana de contraste de 48
  horas, en lugar de una validación cruzada expandida de múltiples pliegues, debido al tiempo de cómputo
  que exige comparar 32 configuraciones por serie. Esta decisión probablemente explica, al menos en parte,
  la baja consistencia entre validación y prueba observada en ALB: una sola ventana de 48 horas es una
  muestra pequeña del comportamiento futuro de la serie y su resultado puede estar dominado por
  particularidades de esas horas específicas.
  \item \textbf{Historial todavía acotado.} El historial efectivo disponible cubre 86 días (2.063
  observaciones horarias), sensiblemente más que el mes del trabajo anterior, pero aún reducido para
  arquitecturas de Deep Learning con una cantidad considerable de parámetros, que pueden requerir un
  historial más extenso para generalizar de forma estable.
  \item \textbf{Modelos con puntaje de validación contaminado.} Cinco de las 32 configuraciones por serie
  utilizan, para su propio ajuste de hiperparámetros o para la elección de sus miembros, la misma ventana
  de 48 horas luego utilizada para clasificarlas. Estos modelos se documentan y se muestran en las tablas y
  figuras de este informe, pero se excluyen de la selección; de haberse permitido, la regla habría
  seleccionado un modelo distinto en ambas series (Tabla~\ref{tab:sensibilidad}), sin que ello implicara una
  mejora real en el conjunto de prueba.
  \item \textbf{Posible anomalía en las últimas horas de store-service.} Las últimas horas registradas de
  la serie store-service (20:00 a 22:00 UTC del 26 de septiembre de 2026) muestran un patrón que se aparta
  del comportamiento de esas mismas horas en los cuatro días anteriores: una caída inusual a las 20:00 UTC,
  seguida de un valor por encima de lo esperado a las 21:00 UTC. No se identificó un patrón equivalente en
  la serie ALB durante esas mismas horas. Esta observación no fue investigada en profundidad y no se
  descarta que corresponda a un artefacto puntual de exportación de datos más que a un cambio real del
  proceso subyacente; se documenta aquí como una limitación abierta y no como un hallazgo verificado.
  \item \textbf{Estacionalidad semanal no aprovechada.} La estacionalidad semanal, débil pero no nula
  (fuerza de 0,206 en ALB y 0,311 en store-service según la descomposición MSTL), se excluyó de los modelos
  Prophet y NeuralProphet y del período estacional del MASE, sobre la base de un historial que, con 86
  días, apenas cubre unas doce semanas completas; un historial más extenso podría justificar su
  incorporación.
  \item \textbf{Serie pública con alcance acotado.} La serie de Wikipedia se evaluó con seis modelos de
  hiperparámetros por defecto, sin Deep Learning, AutoML ni modelos de fundación, con una única ventana de
  validación de 48 horas y una de prueba de 105 horas, y con un SARIMA cuya especificación proviene de ALB.
  Ridge y LightGBM reciben como variable exógena la marca de intervención del despliegue, que vale 1 entre el 17 y el 22 de septiembre aunque esta serie no fue afectada; con la marca en cero, el MASE de prueba de Ridge pasa de 0,784 a 0,756 y el de LightGBM no cambia, sin alterar el orden de los modelos ni el de validación.
  Sus resultados no son comparables uno a uno con los de las 32 configuraciones de las series de BodegaAI.
\end{enumerate}

\subsection{Líneas de trabajo futuro}

La continuación más directa de este trabajo consiste en reemplazar la ventana única de validación por un
esquema de validación cruzada de múltiples orígenes (\emph{multi-origin validation}), en la línea de lo
propuesto por \textcite{hyndman2021}, que estime el error de generalización sobre varias ventanas
sucesivas en lugar de una sola. Una segunda línea consiste en extender el historial disponible una vez que
transcurra más tiempo desde la puesta en producción del sistema, lo que permitiría evaluar con mayor
solidez tanto la estacionalidad semanal como la estabilidad de las arquitecturas de Deep Learning de mayor
capacidad. Finalmente, la serie pública de Wikipedia se trató aquí solo como una línea base acotada; extenderle la
batería completa (Deep Learning, AutoML, modelos de fundación y búsqueda de hiperparámetros) permitiría
contrastar con mayor solidez si las conclusiones alcanzadas para las dos series de infraestructura de
BodegaAI se sostienen en un dominio de datos distinto.


